\documentclass[11pt]{article}
\usepackage[a4paper,margin=1in]{geometry}
\usepackage{amsmath,amssymb,amsthm}
\usepackage{hyperref}

\newcommand{\dd}{\mathrm{d}}
\newcommand{\cD}{\mathcal{D}}
\newcommand{\cH}{\mathcal{H}}
\newcommand{\cQ}{\mathcal{Q}}
\newcommand{\avg}[1]{\langle #1\rangle}
\renewcommand{\Re}{\operatorname{Re}}
\renewcommand{\Im}{\operatorname{Im}}

\newtheorem{theorem}{Theorem}
\newtheorem{lemma}[theorem]{Lemma}
\newtheorem{proposition}[theorem]{Proposition}
\newtheorem{corollary}[theorem]{Corollary}
\theoremstyle{remark}
\newtheorem{remark}[theorem]{Remark}

\title{The quasi-static limit for the linearized Landau--Coulomb operator}
\author{Working note}
\date{17 November 2026}

\begin{document}
\sloppy
\maketitle

\noindent\emph{Summary.} The earlier notes defined the linear heat-flux
closure as the slow-forcing limit of the potential-driven response, and for
the Landau operator assumed that this limit is given by the stationary
kinetic solution. This note proves it.
\begin{itemize}
  \item \textbf{Theorem~\ref{thm:res}}: for every complex frequency $\omega$
    with $\Im\omega\ge0$ in a sector, the resolvent problem
    $(-i\omega+ikv_x+\Lambda)g=1$ has a unique solution with all polynomial
    moments, bounded uniformly in $\omega$.
  \item \textbf{Corollary~\ref{cor:stat}}: at $\omega=0$ this is the
    variational stationary solution of \texttt{landau\_solvability.tex},
    which is therefore unique in a natural class and has all moments.
  \item \textbf{Theorem~\ref{thm:limit}}: as $\omega\to0$ within the sector,
    every polynomial moment of the solution converges to that of the
    stationary solution. Hence $\hat q/\hat T\to q_1/T_1$.
  \item \textbf{Proposition~\ref{prop:time}}: the Laplace transform of any
    weak solution of the time-dependent problem in a weighted energy class
    is given by this resolvent solution.
\end{itemize}
The proofs use the published estimates (C1)--(C5) listed in
\texttt{landau\_solvability.tex}, with the same caveat about sources.
\textbf{Not written out:} existence of the time-dependent weak solution in
the class of Proposition~\ref{prop:time}; it follows by the same
regularization, and is standard in the weighted framework of the cited
papers, but is not proved here. No second reader has checked the proofs.

\section{Setting}

Notation and the estimates (C1)--(C5) are those of
\texttt{landau\_solvability.tex}, except that the velocity weight is written
$m(v)=(1+|v|^2)^{1/2}$, to free the letter $w$. Functions are complex valued
and axisymmetric about the $x$ axis; $\cD$ is extended sesquilinearly, so
$\cD[g]=\cD[\Re g]+\cD[\Im g]$. The spaces $\cH_\theta$ and
$\cH_\infty=\bigcap_\theta\cH_\theta$ are as before. Recall:
\begin{itemize}
  \item $\cD[g]\ge\delta\|g-\Pi g\|^2_\sigma$ and
    $\avg{m^{2\theta-1}|g|^2}\le C_\theta\|g\|^2_{\sigma,\theta}$;
  \item $\Re\cD(g,m^{2\theta}g)\ge\tfrac12\|g\|^2_{\sigma,\theta}
    -C_\theta\int M\rho\,|g|^2$ with $\rho$ a fixed Gaussian;
  \item $|\cD(g_1,m^{2\theta}g_2)|\le C_\theta\|g_1\|_{\sigma,\theta}\|g_2\|_{\sigma,\theta}$.
\end{itemize}
Fix $k\ne0$ and a constant $\Gamma\ge1$, and let
\[
  S_\Gamma=\{\omega\in\mathbb{C}:\ \Im\omega\ge0,\ |\omega|\le1,\
  |\omega|^2\le\Gamma\,\Im\omega\}.
\]
This contains $0$, the segment $\omega=i\varepsilon$, and every non-tangential
approach to $0$ from the upper half-plane.

\paragraph{Weak solutions.} $g\in\cH_1$ is a weak solution of
\begin{equation}
  (-i\omega+ikv_x+\Lambda)\,g=1
  \label{eq:res}
\end{equation}
if $-i\omega\avg{g\bar\varphi}+ik\avg{v_xg\bar\varphi}+\cD(g,\varphi)
=\avg{\bar\varphi}$ for all $\varphi\in\cH_\infty$. All terms are finite: for
$g\in\cH_1$, $\avg{m|g|^2}<\infty$.

\section{The resolvent problem}

\begin{theorem}\label{thm:res}
For every $\omega\in S_\Gamma$ there is a unique weak solution
$g_\omega\in\cH_1$ of \eqref{eq:res}. It lies in $\cH_\infty$, and for every
$\theta\ge0$
\begin{equation}
  \|g_\omega\|_{\sigma,\theta}+(\Im\omega)^{1/2}\avg{m^{2\theta}|g_\omega|^2}^{1/2}
  \le C_\theta ,
  \label{eq:unif}
\end{equation}
with $C_\theta$ depending on $k$, $\Gamma$ and $\theta$ but not on $\omega$.
\end{theorem}

\begin{proof}
\emph{Step 1: uniqueness.} Let $d\in\cH_1$ solve the homogeneous equation.
The weak formulation extends by density and continuity to test functions in
$\cH_1$, so $\varphi=d$ is admissible. The real part gives
$\Im\omega\avg{|d|^2}+\cD[d]=0$, since $ik\avg{v_x|d|^2}$ is imaginary and
finite. So $d\in H$, and then $(-i\omega+ikv_x)d=0$ forces $d=0$.

\emph{Step 2: a regularized problem.} For $\eta\in(0,1]$ consider
\begin{equation}
  (-i\omega+\eta\,m^2+ikv_x+\Lambda)\,g=1 .
  \label{eq:reg}
\end{equation}
Let $V_\theta=\{g:\|g\|^2_{\sigma,\theta}+\avg{m^{2\theta+2}|g|^2}<\infty\}$
and, on $V=V_0$,
\[
  a(g,\varphi)=-i\omega\avg{g\bar\varphi}+\eta\avg{m^2g\bar\varphi}
  +ik\avg{v_xg\bar\varphi}+\cD(g,\varphi).
\]
$a$ is bounded on $V$ because $|v_x|\le m$. It is coercive:
$\Re a(g,g)=\Im\omega\avg{|g|^2}+\eta\avg{m^2|g|^2}+\cD[g]$, and
$\|g\|^2_\sigma\le2\|g-\Pi g\|^2_\sigma+C\avg{|g|^2}$. The Lax--Milgram theorem
gives a unique $g=g_{\omega,\eta}\in V$ with $a(g,\varphi)=\avg{\bar\varphi}$
on $V$.

\emph{Step 3: $g_{\omega,\eta}\in V_\theta$ for every $\theta\ge0$.} On
$V_\theta$ let
$b(g',\varphi)=a(g',m^{2\theta}\varphi)+C_\theta\int M\rho\,g'\bar\varphi$ and
$\ell(\varphi)=\avg{m^{2\theta}\bar\varphi}+C_\theta\int M\rho\,g\,\bar\varphi$.
$b$ is bounded on $V_\theta$, using
$|\avg{v_xm^{2\theta}g'\bar\varphi}|\le\avg{m^{2\theta+1}|g'|^2}^{1/2}
\avg{m^{2\theta+1}|\varphi|^2}^{1/2}$ and (C5). It is coercive, because the
streaming term is imaginary on the diagonal and by (C4):
$\Re b(g',g')\ge\Im\omega\avg{m^{2\theta}|g'|^2}+\eta\avg{m^{2\theta+2}|g'|^2}
+\tfrac12\|g'\|^2_{\sigma,\theta}$. So there is $g'\in V_\theta$ with
$b(g',\varphi)=\ell(\varphi)$. Put $d=g'-g\in V$ and take
$\varphi=m^{-2\theta}d\in V_\theta$:
$a(g',d)+C_\theta\int M\rho\,m^{-2\theta}g'\bar d
=\avg{\bar d}+C_\theta\int M\rho\,m^{-2\theta}g\,\bar d$, while
$a(g,d)=\avg{\bar d}$. Subtracting and taking the real part,
$\eta\avg{m^2|d|^2}\le0$, so $g=g'\in V_\theta$.

\emph{Step 4: energy bound, uniform in $\eta$ and $\omega$.} Testing
\eqref{eq:reg} with $g$ and taking the real part,
\begin{equation}
  \Im\omega\avg{|g|^2}+\eta\avg{m^2|g|^2}+\cD[g]=\Re\avg{g}=:n .
  \label{eq:en}
\end{equation}
Let $\Xi$, $\kappa_{\mathrm F}$, $\chi_*$, $G$, $C_*$ and
$\psi_*=(v_x-\Xi/\kappa_{\mathrm F})/k$ be as in the earlier notes; all are
real and in $\cH_\infty$, $1-kv_x\psi_*=\chi_*$,
$\cD[\psi_*]=1/(k^2\kappa_{\mathrm F})$, and $\avg{\psi_*}=0$. Then
\[
  \avg{g}=\avg{g\chi_*}+k\avg{v_xg\,\psi_*},\qquad
  |\avg{g\chi_*}|=|\cD(g,G)|\le C_*^{1/2}\cD[g]^{1/2}.
\]
Testing \eqref{eq:reg} with $\psi_*$:
$ik\avg{v_xg\psi_*}=i\omega\avg{g\psi_*}-\eta\avg{m^2g\psi_*}-\cD(g,\psi_*)$,
so
\[
  |k\avg{v_xg\psi_*}|\le|\omega|\avg{|g|^2}^{1/2}\avg{\psi_*^2}^{1/2}
  +\eta\avg{m^2|g|^2}^{1/2}\avg{m^2\psi_*^2}^{1/2}
  +\cD[g]^{1/2}(k^2\kappa_{\mathrm F})^{-1/2}.
\]
By \eqref{eq:en}, $\avg{|g|^2}\le n/\Im\omega$,
$\eta\avg{m^2|g|^2}\le n$ and $\cD[g]\le n$. Since
$|\omega|^2\le\Gamma\Im\omega$ and $\eta\le1$,
\[
  n\le|\avg{g}|\le n^{1/2}\Bigl(C_*^{1/2}+(k^2\kappa_{\mathrm F})^{-1/2}
  +\Gamma^{1/2}\avg{\psi_*^2}^{1/2}+\avg{m^2\psi_*^2}^{1/2}\Bigr)
  =:n^{1/2}N_0 .
\]
Hence $n\le N_0^2$ and $|\avg{g}|\le N_0^2$. (For $\omega=0$ the term with
$\avg{|g|^2}$ is absent.)

\emph{Step 5: the component in $H$.} Within the axisymmetric class,
$\Pi g=a+b|v|^2+\beta v_x$. Testing \eqref{eq:reg} with $1$ and with $v_x$,
\[
  ik\avg{v_xg}=1+i\omega\avg{g}-\eta\avg{m^2g},\qquad
  ik\avg{v_x^2g}=i\omega\avg{v_xg}-\eta\avg{m^2v_xg}.
\]
By Step 4, $|\eta\avg{m^2g}|\le\eta\avg{m^2|g|^2}^{1/2}\avg{m^2}^{1/2}\le CN_0$
and similarly for the other term, so $\avg{g}$, $\avg{v_xg}$, $\avg{v_x^2g}$
are bounded. These three numbers, together with $g-\Pi g$, determine
$(a,b,\beta)$ continuously, because $\avg{v_x^2(g-\Pi g)}$ is bounded by
$C\|g-\Pi g\|_\sigma$. Hence $\|g\|_\sigma\le C$ and
$\int M\rho|g|^2\le C$, uniformly.

\emph{Step 6: weighted bound.} By Step 3 the test function $m^{2\theta}g$ is
admissible in \eqref{eq:reg}. The real part, with (C4), gives
\[
  \Im\omega\avg{m^{2\theta}|g|^2}+\eta\avg{m^{2\theta+2}|g|^2}
  +\tfrac12\|g\|^2_{\sigma,\theta}
  \le C_\theta\int M\rho|g|^2+|\avg{m^{2\theta}g}|,
\]
and $|\avg{m^{2\theta}g}|\le\avg{m^{2\theta+1}}^{1/2}\avg{m^{2\theta-1}|g|^2}^{1/2}
\le C\|g\|_{\sigma,\theta}$. With Step 5 this is \eqref{eq:unif} for
$g_{\omega,\eta}$, uniformly in $\eta$.

\emph{Step 7: $\eta\to0$.} By \eqref{eq:unif} a subsequence converges weakly
in every $\cH_\theta$ to some $g_\omega\in\cH_\infty$ satisfying the same
bound. For $\varphi\in\cH_\infty$ each term of the weak form passes to the
limit: $g\mapsto\avg{v_xg\bar\varphi}$ and $g\mapsto\avg{g\bar\varphi}$ are
continuous on $\cH$ because $m^{1/2}v_x\varphi\in L^2(M)$; $\cD(\cdot,\varphi)$
is continuous on $\cH$; and
$\eta|\avg{m^2g\bar\varphi}|\le\eta\,C\to0$. So $g_\omega$ is a weak solution
of \eqref{eq:res}, unique by Step 1.
\end{proof}

\begin{corollary}\label{cor:stat}
Write $g_0=g_e+iw$ for the solution at $\omega=0$. Then $g_e$ is real and
even in $v_x$, $w$ is real and odd, and $g_e$ is the variational solution of
\texttt{landau\_solvability.tex}. Consequently that solution lies in
$\cH_\infty$, and the stationary problem has exactly one weak solution in
$\cH_1$.
\end{corollary}

\begin{proof}
The map $g(v)\mapsto\overline{g(-v_x,v_\perp)}$ sends weak solutions at
$\omega=0$ to weak solutions, so by uniqueness $g_0$ has the stated parity.
Testing with $v_x$ gives $\avg{v_x^2g_e}=0$, so $g_e\in X^\infty$; the
imaginary part of the equation says $f_{g_e}=-w/k$ up to a multiple of
$v_x$. For $\varphi\in X^\infty$ the real part gives
$\cQ(g_e,\varphi)=\cD(g_e,\varphi)-k\avg{v_xw\varphi}=\avg{\varphi}$. Since
$X^\infty$ is dense in $\bar X$, $g_e$ is the unique solution of the
variational problem.
\end{proof}

\section{The limit}

\begin{theorem}\label{thm:limit}
As $\omega\to0$ in $S_\Gamma$, $g_\omega\to g_0$ weakly in $\cH_\theta$ for
every $\theta$, and $\avg{p\,g_\omega}\to\avg{p\,g_0}$ for every polynomial
$p$. In particular, with $Q=\tfrac12v_x(|v|^2-5)$ and
$\zeta=\tfrac13(|v|^2-3)$,
\[
  \avg{Q\,g_\omega}\to q_1=\frac{i}{k},\qquad
  \avg{\zeta\,g_\omega}\to T_1 ,
\]
and whenever $T_1\ne0$, which holds for
$\kappa_{\mathrm F}X_{\max}k^2<3$,
\[
  \frac{\hat q(k,\omega)}{\hat T(k,\omega)}
  =\frac{\avg{Qg_\omega}}{\avg{\zeta g_\omega}}\;\longrightarrow\;
  \frac{q_1}{T_1}=-\,i\,\operatorname{sgn}(k)\,\theta(k).
\]
\end{theorem}

\begin{proof}
Let $\omega_j\to0$ in $S_\Gamma$. By \eqref{eq:unif} a subsequence of
$g_{\omega_j}$ converges weakly in every $\cH_\theta$ to some
$g_*\in\cH_\infty$. In the weak form, $|\omega_j\avg{g_{\omega_j}\bar\varphi}|
\le|\omega_j|(\Im\omega_j)^{-1/2}C\avg{|\varphi|^2}^{1/2}
\le\Gamma^{1/4}|\omega_j|^{1/2}C\to0$, using \eqref{eq:unif} and
$|\omega|\le(\Gamma\Im\omega)^{1/2}$; the other terms converge as in Step 7.
So $g_*$ is a weak solution at $\omega=0$, hence $g_*=g_0$ by
Theorem~\ref{thm:res}. The limit being independent of the subsequence, the
whole family converges. Polynomials define continuous functionals on $\cH$,
since $|\avg{pg}|\le\avg{m\,p^2}^{1/2}\avg{m^{-1}|g|^2}^{1/2}$. The values of
$q_1$ and $T_1$, and the sign of $T_1$ for small $k$, are from
\texttt{landau\_3d\_closure.tex}. The identity
$\hat q/\hat T=\avg{Qg_\omega}/\avg{\zeta g_\omega}$ is
Proposition~\ref{prop:time}.
\end{proof}

\section{Link to the time-dependent problem}

\begin{proposition}\label{prop:time}
Let $a\in L^\infty(0,\infty)$ and let $h(t,v)$ be the amplitude of the
Fourier mode $k$ of a solution of
$\partial_th+ikv_xh+\Lambda h=a(t)\,v_x$, $h(0)=0$, in the following weak
sense: $h\in L^\infty_{\mathrm{loc}}([0,\infty);L^2(m\,M))\cap
L^2_{\mathrm{loc}}([0,\infty);\cH_1)$, with at most polynomial growth in $t$
of these norms, and the equation holds when tested against
$\cH_\infty$-valued smooth functions of $t$. Then for $\Im\omega>0$ the
Laplace transform $\hat h(\omega)=\int_0^\infty e^{i\omega t}h(t)\,\dd t$ is
\[
  \hat h(\omega)=\frac{\hat a(\omega)}{ik}
  +\frac{\omega\,\hat a(\omega)}{k}\,g_\omega ,
\]
and therefore $\hat T=(\omega\hat a/k)\avg{\zeta g_\omega}$,
$\hat q=(\omega\hat a/k)\avg{Qg_\omega}$.
\end{proposition}

\begin{proof}
By the growth assumption $\hat h(\omega)\in\cH_1$ for $\Im\omega>0$, and
integrating the weak form against $e^{i\omega t}$ shows that it is a weak
solution of $(-i\omega+ikv_x+\Lambda)\hat h=\hat a\,v_x$. Since
$(ikv_x+\Lambda)1=ikv_x$, the function
$r=\hat h-\hat a/(ik)$ is a weak solution in $\cH_1$ of
$(-i\omega+ikv_x+\Lambda)r=(\omega\hat a/k)\cdot1$. Uniqueness (Step 1 of
Theorem~\ref{thm:res}, which holds for every $\Im\omega\ge0$) gives
$r=(\omega\hat a/k)g_\omega$; for $\omega$ outside $S_\Gamma$, $g_\omega$ is
defined by the same construction with constants depending on $\omega$.
Constants have $T=q=0$.
\end{proof}

\begin{remark}[What is not written out]
Existence of a time-dependent solution in the class of
Proposition~\ref{prop:time} is obtained by the method of
Theorem~\ref{thm:res}: add $\eta m^2$, solve by Galerkin or semigroup
methods, derive the energy and weighted bounds, which are the time-dependent
versions of Steps 4 and 6, and let $\eta\to0$. Weighted energy estimates of
exactly this type are the basis of the nonlinear theory in the cited papers.
I have not written this argument out, so the statement ``the quasi-static
limit of the time-dependent problem equals the stationary solution'' is
proved here for every weak solution in the stated class, with existence of
such a solution taken from that standard method.
\end{remark}

\section{Summary of status for the Landau operator}

\begin{center}
\renewcommand{\arraystretch}{1.3}
\begin{tabular}{p{8.2cm}p{6.2cm}}
  Statement & Status \\ \hline
  Chapman--Enskog problems solvable, constants finite & proved \\
  Stationary solution exists, unique in $\cH_1$, all moments & proved \\
  Two-sided bound on $\chi(k)$ at small $k$; exact $k^2$ coefficient & proved \\
  Resolvent solution exists, uniform bounds in a sector & proved \\
  Slow-forcing limit of the resolvent equals the stationary solution & proved \\
  Laplace transform of a weak time-dependent solution is the resolvent solution & proved \\
  Existence of the weak time-dependent solution in the weighted class & standard method, not written out \\
  Positivity and $\chi<\kappa_{\mathrm F}$ at all $k$; collisionless limit & open \\
  Tangential approach $\omega\to0$ along the real axis & not covered \\
\end{tabular}
\end{center}
All ``proved'' entries rest on the cited coercivity estimates, one of which
was used through a later restatement.

\end{document}
